% ============================================================
\setpartpreamble{\partintro
  {Wie ein Plugin konfiguriert wird, seine Dienste registriert und sich mit dem DI-Container
  der Anwendung verbindet.}
  {
    \item Kapitel~\ref{ch:the-options-extension}: Die Options-Extension
    \item Kapitel~\ref{ch:registering-services}: Dienste registrieren
    \item Kapitel~\ref{ch:public-api-and-integration-into}: Öffentliche API und Integration
    in die Anwendung
  }}
\part{Infrastruktur}\label{part:II}

% ------------------------------------------------------------
\chapter{Die Options-Extension}\label{ch:the-options-extension}

Der Einstiegspunkt jedes echten Plugins ist eine Klasse, die
\code{IDb\allowbreak{}Context\allowbreak{}Options\allowbreak{}Extension} implementiert. Sie
hat drei Aufgaben: Sie trägt die Konfiguration, sie registriert in
\code{Apply\allowbreak{}Services} die Dienste im internen Provider, und sie beschreibt sich
über \code{Info} gegenüber dem Cache.

\section{Konfiguration als Record}

Alle modellrelevanten Einstellungen des Beispiels liegen in einem Record. Records liefern
\code{Equals} und \code{Get\allowbreak{}Hash\allowbreak{}Code} über alle Felder, was gleich
für den Cache-Schlüssel gebraucht wird. Neue Einstellungen fügt man dem Record hinzu, ohne
Hash und Vergleich anfassen zu müssen.

\begin{sourcecode}[Infrastructure/\allowbreak{}Tenancy\allowbreak{}Settings.\allowbreak{}cs]{csharp}
public sealed record TenancySettings
{
    public static TenancySettings Default { get; } = new();

    public string SettingName { get; init; } = "app.tenant_id";
    public string PolicyName { get; init; } = "tenant_isolation";
    public bool ForceRowLevelSecurity { get; init; } = true;
    public bool ApplyQueryFilter { get; init; }
    public bool CreateTenantIndex { get; init; } = true;
}
\end{sourcecode}

\section{Immutabilität}

EF vergleicht und cached Options. Eine Extension darf sich nach dem Hinzufügen deshalb nie
mehr ändern; jede Änderung erzeugt eine Kopie. Das Muster aus EF Core selbst: ein
geschützter Kopierkonstruktor und eine virtuelle \code{Clone}-Methode, damit abgeleitete
Extensions ihren eigenen Typ klonen.

\begin{sourcecode}[Infrastructure/\allowbreak{}Tenancy\allowbreak{}Options\allowbreak{}Extension.\allowbreak{}cs (\allowbreak{}Auszug)]{csharp}
public class TenancyOptionsExtension : IDbContextOptionsExtension
{
    private DbContextOptionsExtensionInfo? _info;

    public TenancyOptionsExtension()
    {
    }

    protected TenancyOptionsExtension(TenancyOptionsExtension copyFrom)
    {
        ArgumentNullException.ThrowIfNull(copyFrom);

        Settings = copyFrom.Settings;
    }

    public virtual TenancySettings Settings { get; private set; } = TenancySettings.Default;

    public virtual DbContextOptionsExtensionInfo Info => _info ??= new ExtensionInfo(this);

    protected virtual TenancyOptionsExtension Clone() => new(this);

    public virtual TenancyOptionsExtension WithSettings(TenancySettings settings)
    {
        var clone = Clone();
        clone.Settings = settings;
        return clone;
    }
\end{sourcecode}

\section{Die Info-Klasse und der Service-Provider-Cache}

Die Info-Klasse entscheidet, ob EF einen gecachten internen Service Provider wiederverwenden
darf. Dafür gibt es zwei Methoden, die zusammenpassen müssen:
\code{Get\allowbreak{}Service\allowbreak{}Provider\allowbreak{}Hash\allowbreak{}Code} als
schneller Vorfilter und
\code{Should\allowbreak{}Use\allowbreak{}Same\allowbreak{}Service\allowbreak{}Provider} als
exakter Vergleich.

\begin{sourcecode}[Infrastructure/\allowbreak{}Tenancy\allowbreak{}Options\allowbreak{}Extension.\allowbreak{}cs › Extension\allowbreak{}Info]{csharp}
private sealed class ExtensionInfo : DbContextOptionsExtensionInfo
{
    public ExtensionInfo(TenancyOptionsExtension extension)
        : base(extension)
    {
    }

    private new TenancyOptionsExtension Extension => (TenancyOptionsExtension)base.Extension;

    public override bool IsDatabaseProvider => false;

    public override string LogFragment
        => $"AcmeTenancy(Setting={Extension.Settings.SettingName}) ";

    public override int GetServiceProviderHashCode() => Extension.Settings.GetHashCode();

    public override bool ShouldUseSameServiceProvider(DbContextOptionsExtensionInfo other)
        => other is ExtensionInfo o && o.Extension.Settings == Extension.Settings;

    public override void PopulateDebugInfo(IDictionary<string, string> debugInfo)
    {
        debugInfo["Acme:Tenancy:Force"] =
            Extension.Settings.ForceRowLevelSecurity.ToString();
        // ... weitere Einstellungen
    }
}
\end{sourcecode}

Bei diesen beiden Methoden gibt es zwei klassische Fehler mit gegensätzlichen Folgen:

\begin{datatable}{H{0.65}H{1.35}}
  \textbf{Fehler} & \textbf{Folge} \\
  \tablehead
  Modellrelevante Option fehlt im Hash & Zwei Contexts mit unterschiedlicher Konfiguration
  teilen sich Provider und Modell. Der zweite bekommt das gecachte, \textbf{falsche} Modell.
  Fällt im Test oft nicht auf, weil dort nur eine Konfiguration läuft. \\
  \tablerule
  Pro Context variierender Wert im Hash (Mandant, Connection-String, Instanz) & Für jeden
  Context entsteht ein neuer interner Provider samt Modell. Speicher und Startzeit
  explodieren; EF warnt nach 20 Providern mit
  \code{Many\allowbreak{}Service\allowbreak{}Providers\allowbreak{}Created\allowbreak{}Warning}.
  \\
\end{datatable}

\begin{important}[Die Faustregel für den Hash]
  In den Hash gehört genau das, was \textbf{beeinflusst, welche Dienste registriert werden
  oder wie das Modell aussieht}. Alles andere gehört nicht hinein. Der Mandant selbst ändert
  sich pro Request und darf deshalb nie in den Options landen; er kommt zur Laufzeit über
  einen Interceptor (Kapitel~\ref{ch:interceptors}).
\end{important}

\code{Log\allowbreak{}Fragment} erscheint in EFs Startmeldung („initialized using provider …
with options: …“) und hilft bei der Fehlersuche.
\code{Populate\allowbreak{}Debug\allowbreak{}Info} liefert Werte für EFs Diagnose, wenn zu
viele Service Provider entstehen. \code{Is\allowbreak{}Database\allowbreak{}Provider} ist
nur für echte Provider \code{true}.

\section{Validierung}

\code{Validate} wird beim Erstellen des Service Providers aufgerufen. Hier prüft das Plugin
Voraussetzungen, die sich aus der Kombination von Extensions ergeben, und meldet sie mit
einer verständlichen Fehlermeldung statt eines späteren Folgefehlers:

\begin{sourcecode}[Infrastructure/\allowbreak{}Tenancy\allowbreak{}Options\allowbreak{}Extension.\allowbreak{}cs › Validate]{csharp}
public virtual void Validate(IDbContextOptions options)
{
    if (options.FindExtension<NpgsqlOptionsExtension>() is null)
        throw new InvalidOperationException(
            "UseAcmeTenancy() erfordert UseNpgsql(). Für andere Datenbanken wird ein "
            + "eigenes Provider-Paket benötigt.");
}
\end{sourcecode}

% ------------------------------------------------------------
\chapter{Dienste registrieren}\label{ch:registering-services}

\code{Apply\allowbreak{}Services} füllt den internen Service Provider. EF stellt dafür einen
Builder bereit, der die Lebensdauer und Kardinalität jedes EF-Dienstes kennt:
\code{Entity\allowbreak{}Framework\allowbreak{}Services\allowbreak{}Builder} für Kerndienste
und
\code{Entity\allowbreak{}Framework\allowbreak{}Relational\allowbreak{}Services\allowbreak{}Builder}
zusätzlich für relationale Dienste. Er registriert jeden Dienst mit der richtigen
Lebensdauer und nutzt für mehrfach registrierbare Dienste automatisch
\code{Try\allowbreak{}Add\allowbreak{}Enumerable}.

\begin{sourcecode}[Infrastructure/\allowbreak{}Tenancy\allowbreak{}Options\allowbreak{}Extension.\allowbreak{}cs › Apply\allowbreak{}Services]{csharp}
public virtual void ApplyServices(IServiceCollection services)
{
    // 1) Plugin-Interfaces: mehrfach registrierbar, kombinierbar mit anderen Plugins
    new EntityFrameworkRelationalServicesBuilder(services)
        .TryAdd<IRelationalTypeMappingSourcePlugin, TenancyTypeMappingSourcePlugin>()
        .TryAdd<IMethodCallTranslatorPlugin, TenancyMethodCallTranslatorPlugin>()
        .TryAdd<IMemberTranslatorPlugin, TenancyMemberTranslatorPlugin>()
        .TryAdd<IEvaluatableExpressionFilterPlugin,
                TenancyEvaluatableExpressionFilterPlugin>()
        .TryAdd<IConventionSetPlugin, TenancyConventionSetPlugin>()
        .TryAdd<ISingletonOptions, ITenancySingletonOptions>(
            p => p.GetRequiredService<ITenancySingletonOptions>())
        // 2) Eigene Dienste, die EF nicht kennt
        .TryAddProviderSpecificServices(m => m
            .TryAddSingleton<ITenancySingletonOptions, TenancySingletonOptions>()
            .TryAddSingleton<TenancyLoggingDefinitions, TenancyLoggingDefinitions>());

    // 3) Einzeldienste des Providers ersetzen
    services.Replace(ServiceDescriptor
        .Scoped<IMigrationsSqlGenerator, TenancyMigrationsSqlGenerator>());
    services.Replace(ServiceDescriptor
        .Singleton<IRelationalAnnotationProvider, TenancyAnnotationProvider>());
}
\end{sourcecode}

\begin{datatable}{H{0.83}H{0.77}H{1.40}}
  \textbf{Methode} & \textbf{Wann} & \textbf{Warum} \\
  \tablehead
  \code{Try\allowbreak{}Add<\allowbreak{}TService,\allowbreak{} TImpl>(\allowbreak{})} &
  EF-Dienste und Plugin-Interfaces & Lebensdauer und Mehrfachregistrierung kommen aus EFs
  Tabelle. Falsche Lebensdauern führen sonst zu Exceptions oder subtil falschem Verhalten.
  \\
  \tablerule
  \code{Try\allowbreak{}Add\allowbreak{}Provider\allowbreak{}Specific\allowbreak{}Services}
  & Eigene Dienste, die EF nicht kennt & Der Builder lehnt unbekannte Typen ab; hier legt
  man die Lebensdauer selbst fest. \\
  \tablerule
  \code{services.\allowbreak{}Replace(\allowbreak{}.\allowbreak{}.\allowbreak{}.\allowbreak{})}
  & Einzeldienst des Providers ersetzen & Der Provider registriert per
  \code{Try\allowbreak{}Add}. \code{Replace} gewinnt unabhängig davon, ob das Plugin vor
  oder nach dem Provider läuft. \\
\end{datatable}

\begin{note}[Warum Replace in beiden Reihenfolgen funktioniert]
  EF ruft \code{Apply\allowbreak{}Services} der Extensions in der Reihenfolge auf, in der
  sie hinzugefügt wurden. Läuft das Plugin zuerst, fügt \code{Replace} den Dienst hinzu und
  das spätere \code{Try\allowbreak{}Add} des Providers greift nicht mehr. Läuft es danach,
  ersetzt \code{Replace} den Provider-Dienst. Beide Reihenfolgen führen zum selben Ergebnis.
\end{note}

\section{Singleton-Options}

Singleton-Dienste wie Type-Mapping-Plugins können
\code{IDb\allowbreak{}Context\allowbreak{}Options} nicht pro Context auflösen. Für sie gibt
es \code{ISingleton\allowbreak{}Options}: ein Singleton, das EF einmal beim Erstellen des
Providers mit den Options initialisiert und bei jedem weiteren Context validiert.

\begin{sourcecode}[Infrastructure/\allowbreak{}ITenancy\allowbreak{}Singleton\allowbreak{}Options.\allowbreak{}cs]{csharp}
public class TenancySingletonOptions : ITenancySingletonOptions
{
    public TenancySettings Settings { get; private set; } = TenancySettings.Default;

    // Einmal pro internem Service Provider
    public virtual void Initialize(IDbContextOptions options)
        => Settings = options.FindExtension<TenancyOptionsExtension>()?.Settings
                      ?? TenancySettings.Default;

    // Bei jedem Context: Optionen dürfen nicht vom Stand des Providers abweichen
    public virtual void Validate(IDbContextOptions options)
    {
        var settings = options.FindExtension<TenancyOptionsExtension>()?.Settings
                       ?? TenancySettings.Default;
        if (settings != Settings)
            throw new InvalidOperationException("Tenancy-Optionen nachträglich geändert ...");
    }
}
\end{sourcecode}

Die doppelte Registrierung in \code{Apply\allowbreak{}Services} ist nötig: einmal unter dem
eigenen Interface, damit Dienste es injizieren können, und einmal als
\code{ISingleton\allowbreak{}Options}, damit EF es findet und initialisiert. Die
Factory-Variante sorgt dafür, dass beide Registrierungen \emph{dieselbe} Instanz liefern.

\section{Kombinierbarkeit trotz Ersetzen}

Ersetzt ein Plugin einen Einzeldienst, kann ein zweites Plugin, das dasselbe tut, nicht
gleichzeitig aktiv sein. Professionelle Plugins entschärfen das, indem sie die eigentliche
Logik zusätzlich öffentlich und wiederverwendbar anbieten. Im Beispiel ist
\code{Create\allowbreak{}Policy\allowbreak{}Sql} öffentlich und alle Generator-Methoden sind
virtuell. Ein Nutzer, der zwei Plugins kombinieren muss, schreibt einen eigenen Generator,
der von einem erbt und die Logik des anderen aufruft. Diese Möglichkeit gehört in die
Dokumentation.

Noch besser vermeidet ein Plugin den Konflikt ganz, indem es den Dienst \textbf{dekoriert}
statt ersetzt: Es nimmt die Implementierung, die zu diesem Zeitpunkt registriert ist (die
des Providers oder eines anderen Plugins), umhüllt sie und delegiert an sie. Ein Beispiel
ist ein (fiktives) Audit-Plugin, das an jede neue Tabelle einen Trigger hängt. In seinem
\code{Apply\allowbreak{}Services} sucht es die vorhandene Registrierung und ersetzt nur den
Deskriptor, nicht die Logik:

\begin{sourcecode}[tests/\allowbreak{}Acme.\allowbreak{}Entity\allowbreak{}Framework\allowbreak{}Core.\allowbreak{}Tenancy.\allowbreak{}Tests/\allowbreak{}Composition\allowbreak{}Tests.\allowbreak{}cs › Audit\allowbreak{}Options\allowbreak{}Extension]{csharp}
public void ApplyServices(IServiceCollection services)
{
    // Bisher registriert: der Generator des Providers oder eines anderen Plugins
    var inner = services.LastOrDefault(d => d.ServiceType == typeof(IMigrationsSqlGenerator))
        ?? throw new InvalidOperationException("Call UseContosoAudit() after UseNpgsql().");

    services.Replace(ServiceDescriptor.Describe(
        typeof(IMigrationsSqlGenerator),
        sp => new AuditMigrationsSqlGenerator(
            (IMigrationsSqlGenerator)CreateInstance(sp, inner),
            sp.GetRequiredService<MigrationsSqlGeneratorDependencies>()),
        inner.Lifetime));     // Lebensdauer des Originals
}

// Erzeugt die Implementierung so, wie es der ursprüngliche Deskriptor getan hätte
private static object CreateInstance(IServiceProvider sp, ServiceDescriptor descriptor)
    => descriptor.ImplementationInstance
       ?? descriptor.ImplementationFactory?.Invoke(sp)
       ?? ActivatorUtilities.CreateInstance(sp, descriptor.ImplementationType!);
\end{sourcecode}

Der Trick funktioniert, weil EF \code{Apply\allowbreak{}Services} in der Reihenfolge
aufruft, in der die Extensions hinzugefügt wurden; wenn das Audit-Plugin läuft, hat der
Provider seinen Generator bereits registriert. \code{Create\allowbreak{}Instance} bildet die
ursprüngliche Registrierung exakt nach: eine Instanz, eine Factory oder einen Typ, dessen
Abhängigkeiten (etwa Npgsqls Singleton-Options) \code{Activator\allowbreak{}Utilities} aus
dem Container auflöst. Der Decorator selbst nutzt nur die öffentliche Schnittstelle:

\begin{sourcecode}[tests/\allowbreak{}.\allowbreak{}.\allowbreak{}.\allowbreak{}/\allowbreak{}Composition\allowbreak{}Tests.\allowbreak{}cs › Audit\allowbreak{}Migrations\allowbreak{}Sql\allowbreak{}Generator (\allowbreak{}gekürzt)]{csharp}
public class AuditMigrationsSqlGenerator : IMigrationsSqlGenerator
{
    // Konstruktor: (IMigrationsSqlGenerator inner, MigrationsSqlGeneratorDependencies deps)

    public IReadOnlyList<MigrationCommand> Generate(
        IReadOnlyList<MigrationOperation> operations,
        IModel? model = null,
        MigrationsSqlGenerationOptions options = MigrationsSqlGenerationOptions.Default)
    {
        var commands = _inner.Generate(operations, model, options).ToList();   // unverändert
        var builder = new MigrationCommandListBuilder(_dependencies);
        var sql = _dependencies.SqlGenerationHelper;

        // Am Ende angehängt: dann existieren alle Tabellen der Migration
        foreach (var table in operations.OfType<CreateTableOperation>())
        {
            builder.Append("CREATE TRIGGER audit_changes ")
                .Append("AFTER INSERT OR UPDATE OR DELETE ON ")
                .Append(sql.DelimitIdentifier(table.Name, table.Schema))
                .Append(" FOR EACH ROW EXECUTE FUNCTION audit.log_change()")
                .AppendLine(sql.StatementTerminator)
                .EndCommand();
        }

        commands.AddRange(builder.GetCommandList());
        return commands;
    }
}
\end{sourcecode}

In der Anwendung lassen sich beide Plugins ohne eine Zeile Verbindungscode kombinieren:

\begin{sourcecode}[Program.\allowbreak{}cs]{csharp}
services.AddDbContext<ShopContext>((sp, options) => options
    .UseNpgsql(connectionString)
    .UseAcmeTenancy(sp)       // ersetzt IMigrationsSqlGenerator
    .UseContosoAudit());      // umhüllt, was registriert ist: hier Acmes
\end{sourcecode}

\begin{important}[Dekorieren, wo möglich; ersetzen, wo nötig]
  Decorators lassen sich in beliebiger Reihenfolge kombinieren, weil jeder umhüllt, was er
  vorfindet. Ein Plugin, das \code{Replace} verwendet, verwirft dagegen alles, was vor ihm
  registriert wurde, auch Decorators. Ein ersetzendes Plugin muss deshalb zuerst aktiviert
  werden, und ein Decorator muss prüfen, ob es etwas zu dekorieren gibt, und sonst mit einer
  klaren Meldung abbrechen (Provider fehlt oder wurde später hinzugefügt).

  Dekorieren setzt voraus, dass die öffentliche Schnittstelle genügt. Das Tenancy-Plugin
  braucht geschützte Einstiegspunkte innerhalb von \code{CREATE TABLE} und erbt deshalb vom
  Npgsql-Generator. Für solche Fälle bleibt der oben beschriebene Ausweg: öffentliche Logik
  wie \code{Create\allowbreak{}Policy\allowbreak{}Sql}, virtuelle Methoden und
  \code{Replace\allowbreak{}Service} in der Anwendung.
\end{important}

\begin{tip}
  Der Test \code{Composition\allowbreak{}Tests} im Beispiel enthält das vollständige
  Audit-Plugin. Er prüft, dass der Decorator Acmes Generator umhüllt, dass die Migration das
  SQL beider Plugins enthält und dass das Aktivieren vor dem Provider fehlschlägt.
\end{tip}

% ------------------------------------------------------------
\chapter{Öffentliche API und Integration in die Anwendung}\label{ch:public-api-and-integration-into}

\section{Die Use-Methode}

Nutzer aktivieren das Plugin über eine Extension-Method auf
\code{Db\allowbreak{}Context\allowbreak{}Options\allowbreak{}Builder}. Sie holt die
vorhandene Extension oder legt eine neue an und registriert sie über
\code{IDb\allowbreak{}Context\allowbreak{}Options\allowbreak{}Builder\allowbreak{}Infrastructure}:

\begin{sourcecode}[Infrastructure/\allowbreak{}Tenancy\allowbreak{}Db\allowbreak{}Context\allowbreak{}Options\allowbreak{}Builder\allowbreak{}Extensions.\allowbreak{}cs (\allowbreak{}Auszug)]{csharp}
public static DbContextOptionsBuilder UseAcmeTenancy(
    this DbContextOptionsBuilder optionsBuilder,
    Action<TenancyOptionsBuilder>? configure = null)
{
    ArgumentNullException.ThrowIfNull(optionsBuilder);

    var extension = optionsBuilder.Options.FindExtension<TenancyOptionsExtension>()
                    ?? new TenancyOptionsExtension();

    ((IDbContextOptionsBuilderInfrastructure)optionsBuilder).AddOrUpdateExtension(extension);
    configure?.Invoke(new TenancyOptionsBuilder(optionsBuilder));
    return optionsBuilder;
}

// Generische Überladung: Method Chaining bleibt typisiert
public static DbContextOptionsBuilder<TContext> UseAcmeTenancy<TContext>(
    this DbContextOptionsBuilder<TContext> optionsBuilder,
    Action<TenancyOptionsBuilder>? configure = null) where TContext : DbContext
    => (DbContextOptionsBuilder<TContext>)UseAcmeTenancy(
        (DbContextOptionsBuilder)optionsBuilder, configure);
\end{sourcecode}

Die Datei liegt im Namespace
\code{Microsoft.\allowbreak{}Entity\allowbreak{}Framework\allowbreak{}Core}. So ist die
Methode ohne zusätzliches \code{using} auffindbar, genau wie \code{Use\allowbreak{}Npgsql}.
Das ist die verbreitete Konvention im EF-Ökosystem.

\section{Der Options-Builder}

Einstellungen werden über einen eigenen Builder gesetzt. Er validiert Eingaben sofort und
schreibt jede Änderung als neue Extension zurück:

\begin{sourcecode}[Infrastructure/\allowbreak{}Tenancy\allowbreak{}Options\allowbreak{}Builder.\allowbreak{}cs (\allowbreak{}Auszug)]{csharp}
public class TenancyOptionsBuilder
{
    public TenancyOptionsBuilder(DbContextOptionsBuilder optionsBuilder)
    {
        ArgumentNullException.ThrowIfNull(optionsBuilder);

        OptionsBuilder = optionsBuilder;
    }

    protected virtual DbContextOptionsBuilder OptionsBuilder { get; }

    public virtual TenancyOptionsBuilder UseSettingName(string settingName)
        => With(s => s with
        {
            SettingName = TenancyIdentifiers.ValidateSettingName(settingName),
        });

    public virtual TenancyOptionsBuilder ApplyQueryFilter(bool apply = true)
        => With(s => s with { ApplyQueryFilter = apply });

    protected virtual TenancyOptionsBuilder With(
        Func<TenancySettings, TenancySettings> change)
    {
        var extension = OptionsBuilder.Options.FindExtension<TenancyOptionsExtension>()
                        ?? new TenancyOptionsExtension();

        ((IDbContextOptionsBuilderInfrastructure)OptionsBuilder)
            .AddOrUpdateExtension(extension.WithSettings(change(extension.Settings)));
        return this;
    }
}
\end{sourcecode}

\begin{note}[Alternative: Aktivierung über den Provider-Builder]
  Provider-spezifische Plugins werden manchmal auch über den Options-Builder des Providers
  aktiviert, etwa \code{Use\allowbreak{}Npgsql(\allowbreak{}cs,\allowbreak{} o =>
  o.\allowbreak{}Use\allowbreak{}Noda\allowbreak{}Time(\allowbreak{}))}. Dafür schreibt man
  eine Extension-Method auf
  \code{Npgsql\allowbreak{}Db\allowbreak{}Context\allowbreak{}Options\allowbreak{}Builder}
  und kommt über
  \code{IRelational\allowbreak{}Db\allowbreak{}Context\allowbreak{}Options\allowbreak{}Builder\allowbreak{}Infrastructure.\allowbreak{}Options\allowbreak{}Builder}
  an den eigentlichen Builder. Das Beispiel nutzt den einfacheren Weg über
  \code{Db\allowbreak{}Context\allowbreak{}Options\allowbreak{}Builder}, weil das Plugin
  mehr ist als eine Provider-Option.
\end{note}

\section{Zwei Container verbinden}

Manche Teile eines Plugins brauchen Dienste der \emph{Anwendung}: Der Connection-Interceptor
braucht den aktuellen Mandanten, den nur die Anwendung kennt. Der interne EF-Provider kann
App-Dienste aber nicht auflösen. Die saubere Lösung: Solche Komponenten leben im
App-Container und werden per \code{Add\allowbreak{}Interceptors} an EF übergeben.

\begin{sourcecode}[Dependency\allowbreak{}Injection/\allowbreak{}… (\allowbreak{}Kern),\allowbreak{} Infrastructure/\allowbreak{}… (\allowbreak{}Postgre\allowbreak{}SQL),\allowbreak{} Auszüge]{csharp}
// Kernpaket: nur den Mandanten-Provider im App-Container registrieren
public static IServiceCollection AddAcmeTenancy<TTenantProvider>(
    this IServiceCollection services,
    ServiceLifetime lifetime = ServiceLifetime.Scoped)
    where TTenantProvider : class, ITenantProvider
{
    services.TryAdd(ServiceDescriptor.Describe(
        typeof(ITenantProvider), typeof(TTenantProvider), lifetime));
    return services;
}

// Provider-Paket: Use-Überladung erzeugt die Interceptors aus diesem Provider
public static DbContextOptionsBuilder UseAcmeTenancy(
    this DbContextOptionsBuilder optionsBuilder,
    IServiceProvider applicationServices,
    Action<TenancyOptionsBuilder>? configure = null)
{
    var tenants = applicationServices.GetRequiredService<ITenantProvider>();

    return optionsBuilder
        .UseAcmeTenancy(configure)
        .AddInterceptors(
            new TenantConnectionInterceptor(tenants),
            new TenantSaveChangesInterceptor(tenants));
}
\end{sourcecode}

Bei \code{Add\allowbreak{}Db\allowbreak{}Context(\allowbreak{}(\allowbreak{}sp,\allowbreak{}
o) => .\allowbreak{}.\allowbreak{}.\allowbreak{})} ist \code{sp} der Scope der aktuellen
Anfrage, weil EF die Options standardmäßig pro Scope baut. Der Mandanten-Provider und damit
die Interceptors sehen den richtigen Mandanten. Die Registrierung liegt im Kernpaket, weil
sie für jeden Provider gleich ist (Kapitel~\ref{ch:what-belongs-in-the-core}). Für die
Brücke in die andere Richtung kann ein Interceptor über
\code{event\allowbreak{}Data.\allowbreak{}Context.\allowbreak{}Get\allowbreak{}Service<\allowbreak{}T>(\allowbreak{})}
Dienste aus dem \emph{internen} Provider lesen; der Connection-Interceptor holt sich so den
Setting-Namen aus den Singleton-Options (Kapitel~\ref{ch:interceptors}).

\begin{important}[DbContext-Pooling]
  Bei \code{Add\allowbreak{}Db\allowbreak{}Context\allowbreak{}Pool} werden Contexts
  wiederverwendet und bekommen ihre Options nur einmal. Scoped-Interceptors aus dem ersten
  Scope würden für alle späteren Anfragen weiterverwendet, mit dem Mandanten der ersten
  Anfrage. Für Pooling muss der Interceptor ein Singleton sein, der den Mandanten bei jedem
  Aufruf neu ermittelt, etwa über \code{IHttp\allowbreak{}Context\allowbreak{}Accessor}.
  Kapitel~\ref{ch:using-the-plugin-in-asp} zeigt diese Variante. Dasselbe gilt für
  \code{Add\allowbreak{}Db\allowbreak{}Context\allowbreak{}Factory}, dessen Factory
  standardmäßig ein Singleton ist.
\end{important}
